\documentclass[12pt]{article}
\usepackage[utf8]{inputenc}
\usepackage[T1]{fontenc}
\usepackage{hyperref}
\title{Yukuben Language Resources for K-12 Curriculum Development}
\author{Compiled by Claude Sonnet 5.5 (claude-sonnet-5-5) for Andr\'as Kornai}
\date{}
\begin{document}
\maketitle

\section{Introduction}
Yukuben (ISO 639-3 ybl; also Fete, Uhumghikgi/Uhum, Nyikuben) is spoken in Taraba State, Nigeria, and across the border in North-West Cameroon. Wikipedia cites 16{,}000 speakers in Nigeria (1992) and about 1{,}000 in Cameroon, and says it works as a local trade language. Other sources give different figures: Global Recordings Network about 15{,}000, and Joshua Project about 37{,}000 across two people groups. Classification is disputed. Glottolog places it in Plateau, Ethnologue places it in Jukunoid, and Blench (2026) argues for an independent Benue--Congo branch. Endangeredlanguages.com classes it as vulnerable. Educational materials are very scarce in my searches.

\section{Primary Educational Resources}
No primers, textbooks or school materials in Yukuben turned up in my searches.

\section{Linguistic and Reference Materials}
\begin{itemize}
\item Wikipedia, Yukuben language: \url{https://en.wikipedia.org/wiki/Yukuben_language}
\item Wikipedia, Yukubenic languages: \url{https://en.wikipedia.org/wiki/Yukubenic_languages}
\item Glottolog entry (yuku1243): \url{https://glottolog.org/resource/languoid/id/yuku1243}
\item Endangered Languages Project entry: \url{https://endangeredlanguages.com/lang/6107}
\end{itemize}
These are overview pages. No dictionary or grammar turned up in my searches.

\section{Religious and Community Resources}
Joshua Project lists Bible translation as started, with no audio Bible or Jesus Film available.
\begin{itemize}
\item Joshua Project, Uhumkighi/Yukuben in Cameroon: \url{https://www.joshuaproject.net/index.php/people_groups/14097/CM}
\item Joshua Project language page: \url{https://www.joshuaproject.net/index.php/languages/ybl}
\end{itemize}

\section{Digital and Audio Resources}
\begin{itemize}
\item Global Recordings Network, Yukuben audio recordings (evangelism recordings): \url{https://globalrecordings.net/en/language/ybl}
\end{itemize}

\section{Curriculum Development Context}
There is no evidence of Yukuben in schools or of a standard orthography. Hausa and English are the likely school languages in the area, but I did not confirm this in my searches. A curriculum would need new documentation work, and the Global Recordings audio could serve as a seed for listening materials.
\end{document}